% !TeX root = ../tesis.tex
%%%%%%%%%%%%%%%%%%%%%%%%%%%%%%%%%%%%%%%%%%%%%%%%%%%%%%%%%%%%%%%%%%%
% CAPÍTULO 7 — CONCLUSIONES (Issue #23)
%
% Plan y decisiones: 7-Conclusiones/CONTEXTO_CAP7.md (D7-1…D7-4).
% Hilo conductor: «lo que se mide define lo que se construye»
% (Programa Maestro del Metro 1985, §6.5.2).
% Cifras: corrida 2026-10-04 (VFTModel 4ef4265).
%
% Estructura: entrada (resumen + contraste) · 7.1 Tecnológico ·
%             7.2 Urbano · 7.3 Académico y Trabajo Futuro (con cierre)
%%%%%%%%%%%%%%%%%%%%%%%%%%%%%%%%%%%%%%%%%%%%%%%%%%%%%%%%%%%%%%%%%%%
\chapter{Conclusiones}
\label{ch:conclusiones}

% ENTRADA — P1 resumen (✅ 2026-11-04, score 0.122).
La tesis empezó por la forma de la red. En la \zmvm{} el transporte público
creció hacia el centro y dejó a la periferia con servicios que casi no se
cruzan. El Periférico rodea la ciudad, pero se hizo para el automóvil y por él
no corre transporte masivo (Capítulos~\ref{ch:introduccion}
y~\ref{ch:marco-teorico}). Medir eso pedía una herramienta propia: una cadena
de tres capas de código abierto (Capítulo~\ref{ch:Capas-codigo}), ocho
indicadores y la Clasificación Garibelt para leerlos juntos
(Capítulo~\ref{ch:definiciones}). Con esa cadena se diagnosticó la red actual,
y el Espectro MJG quedó como línea base (Capítulo~\ref{ch:analisis-red-actual}).
Luego se probaron sobre el Anillo Periférico Interior dos casos con el mismo
trazo y las mismas 94 estaciones, el Anillo MB y el Anillo Metro, que solo
cambian en la operación (Capítulo~\ref{ch:analisis-red-anillo}).

% P2 — contraste con la hipótesis y veredicto de la pregunta
% (✅ 2026-11-04, score 0.122). DI en la periferia: tendencia sobre
% pocas rutas (tab:di-alcaldia, n = 1–2). Cifras de la pregunta en 7.2 (D7-8).
Frente a la hipótesis (§\ref{sec:hipotesis}), el resultado es una
confirmación parcial. La primera parte se sostiene: la red radial-concéntrica
tiene vulnerabilidades que se pueden medir. La intermediación se concentra en
unos cuantos nodos del núcleo, con Tacubaya a la cabeza, y los rodeos más
largos aparecen en la periferia sur y oriente, aunque sobre pocas rutas
(§\ref{subsec:factor-desviacion}). Los 122 minutos que la \textsc{semovi}
reporta para el Cuarto Contorno \citep{semovi2020diagnostico} no se pueden
contrastar con el modelo, porque su $T$ de 113.91 minutos promedia toda la red.
La segunda parte también se cumple: la metodología midió esas vulnerabilidades
y repitió el cálculo en los tres casos con el mismo código. La tercera se
cumple a medias. Con el anillo bajan el tiempo de viaje y la intermediación de
Tacubaya, pero la concentración no se reparte, y el $DI$ quedó fuera de la
comparación porque solo se calculó para la red actual. La pregunta de
investigación recibe una respuesta del mismo tipo. El anillo mejora la conexión
de la periferia que ya está cerca de la red, no la de la que queda lejos, y el
modelo anticipa la estructura de ese escenario, pero no su demanda
(§\ref{sec:conclusiones-urbano}).

% P3 — objetivos específicos y transición a los tres ejes
% (✅ 2026-11-04, opción A, score 0.122). ΔE del obj. 3 solo
% exploratorio (#39). El pasaje de 1985 lo abre 7.1 y lo retoma el cierre.
De los cinco objetivos específicos (§\ref{subsec:obj-especificos}), tres se
cumplieron completos y dos en parte. Los \termino{feeds} \gtfs{} se
sistematizaron en un grafo multimodal, sobre él se calcularon el $DI$, $T$, $C$
y $FC$ de la red actual, y con la cobertura y la Fuerza Capilar se ubicó el
mayor déficit, sobre todo en la periferia mexiquense. Los nodos críticos
quedaron localizados con la intermediación, pero la caída de eficiencia
$\Delta E$ se quedó en un ejercicio exploratorio. Los dos casos del anillo se
simularon con $C$, $FC$, $T$ y $B(v)$, y el Anillo Metro resultó el más
favorable; faltó el $DI$ para completar la comparación. Las secciones que
siguen ordenan estos resultados en tres ejes. El tecnológico trata del
instrumento que hizo posible medir la red; el urbano, de lo que esa medición
dice sobre ella; el académico, de lo que todavía no mide y del trabajo que de
ahí se desprende.
% Pregunta de investigación: solo el veredicto (confirmación parcial)
% con remisión a §\ref{sec:conclusiones-urbano}; las cifras están ahí.

% !TeX root = ../tesis.tex
%%%%%%%%%%%%%%%%%%%%%%%%%%%%%%%%%%%%%%%%%%%%%%%%%%%%%%%%%%%%%%%%%%%
% SECCIÓN 7.1 — EJE TECNOLÓGICO (Issue #23)
% Papel en el hilo: el instrumento que hace medible la estructura.
% Esquema: problema (1985) → qué se hizo → replicabilidad → aporte →
% límites y transición a 7.2. Plan: CONTEXTO_CAP7.md §4.
% Fuera a propósito: T_b fijo (#32), grafo simple (D9), nodo
% «Memorial», 13 nodos del snapping; TransCAD/NAVTEQ sin verificar.
%%%%%%%%%%%%%%%%%%%%%%%%%%%%%%%%%%%%%%%%%%%%%%%%%%%%%%%%%%%%%%%%%%%
\section{Eje Tecnológico: un Instrumento para Medir la Estructura}
\label{sec:conclusiones-tecnologico}

El Programa Maestro del Metro de 1985 descartó los dos corredores del
Periférico porque no alcanzaban la captación que pedía su red
(§\ref{subsubsec:anillo-pm-2010}). Sus parámetros salían de la demanda, y con
ellos la evaluación no tenía forma de medir lo que el anillo aportaba a la
estructura de la red. Aquel estudio se hizo con otras tecnologías, las de su
momento. Medir la estructura pide hoy otro tipo de herramienta, y por eso el
primer resultado de esta tesis es un instrumento: una cadena abierta que mide
la estructura de la red de transporte con datos públicos.

La cadena tiene tres capas, y cada una cumple una tarea
(Figura~\ref{fig:macro-arquitectura-stack}). \texttt{Apímetro} reúne la
información de doce sistemas de transporte y 141 demarcaciones, que cada
organismo publica a su manera, y la entrega en un solo formato. Con ello cubre
algo que la Tarjeta MI dejó pendiente en 2019, cuando integró la tarifa pero no
los datos (§\ref{subsubsec:espectro-cierre}). \texttt{VFTModel} arma con esa
información un grafo de 11,115 nodos y 45,517 aristas, asigna a cada arista un
tiempo con la ecuación VFT y calcula los ocho indicadores
(§\ref{sec:vft-indicadores}). \texttt{Transport-gis-zmvm-mjg}, por último,
lleva los resultados a mapas de QGIS y a \termino{dashboards} de Tableau, que
es donde un urbanista puede leerlos.

Que la cadena sirva como instrumento, y no solo para resolver un caso, se
comprobó en el Capítulo~\ref{ch:analisis-red-anillo}. Los tres escenarios se
procesaron con el mismo código. Para sumar el anillo bastó registrarlo en
\texttt{Apímetro} como un sistema más, y cada caso corrió en sus propios
contenedores, sin que los datos de uno alteraran los de otro
(§\ref{subsec:anillo-procesamiento}). Con esa separación se pueden comparar
corredores y estimar el efecto de una línea nueva antes de construirla
(§\ref{sec:vft-salidas}). Es el mismo ejercicio que hizo el Programa Maestro al
simular su red varias veces, solo que ahora lo que se compara es la estructura.

El aporte de este eje no se limita a la \zmvm{}. La cadena no depende de
licencias comerciales y se puede repetir en cualquier ciudad que tenga
\termino{feeds} \gtfs{} y cartografía de base, algo que ya cumple la mayoría de las ciudades
latinoamericanas de tamaño parecido \citep{barrington2017world}. El código se
publica bajo la \textit{Business Source License~1.1}, que pasa a
\textit{Apache~2.0} en abril de 2029, y viene con una wiki, un foro,
\texttt{seed.sql} y Docker Compose para que lo use también quien no programa
(§\ref{subsec:futuro-apimetro}). \citet{gurstein2011open} advierte que los
datos abiertos solo empoderan a quien puede procesarlos; una herramienta así
acorta esa distancia. En ese sentido entiende esta tesis la soberanía
tecnológica: una capacidad de análisis que se queda en la institución pública y
no se compra a un proveedor.

El instrumento tiene límites, y conviene tenerlos presentes al leer sus
resultados. Las velocidades de superficie vienen de \texttt{Apímetro} y no se
han contrastado con velocidades observadas por ruta
(§\ref{subsec:validacion-impedancia}). La caída de eficiencia $\Delta E$ se
calculó solo como ejercicio exploratorio, con una corrida anterior a la vigente
(§\ref{subsec:vulnerabilidad-estructural}). Las cinco estaciones del Tren
Interurbano quedan fuera de la componente gigante. El límite principal es otro:
el modelo mide la estructura, no el uso. No calcula demanda ni capacidad, aunque
\texttt{Apímetro} ya guarda la capacidad de los vehículos. Estos pendientes se
retoman como trabajo futuro en §\ref{sec:conclusiones-academico}. Con el
instrumento probado, falta ver qué dice esa medición sobre la red, que es el
tema de §\ref{sec:conclusiones-urbano}.

% !TeX root = ../tesis.tex
%%%%%%%%%%%%%%%%%%%%%%%%%%%%%%%%%%%%%%%%%%%%%%%%%%%%%%%%%%%%%%%%%%%
% SECCIÓN 7.2 — EJE URBANO (Issue #23)
% Papel en el hilo: lo que la medición de la estructura revela.
% Esquema: diagnóstico → lectura social → acercar no es conectar →
% radial a anillo-radial y modo → CETRAM y lectura histórica →
% aporte, pendientes y transición a 7.3. Plan: CONTEXTO_CAP7.md §5.
% Las cifras de la pregunta de investigación quedan aquí; la entrada
% solo da el veredicto y remite a esta sección (decisión 2026-11-03).
% Fuera a propósito: Moscú/Melbourne, $408 M (1958), 92.7 % de
% superficie; «dos ejes, red menos frágil» solo como pendiente.
%%%%%%%%%%%%%%%%%%%%%%%%%%%%%%%%%%%%%%%%%%%%%%%%%%%%%%%%%%%%%%%%%%%
\section{Eje Urbano}
\label{sec:conclusiones-urbano}

Medida la estructura, la red de la \zmvm{} muestra un patrón que la captación
por corredor no registraba: eficiencia interna y exclusión perimetral
(§\ref{subsubsec:espectro-cierre}). La red produce viajes eficientes para quien
ya forma parte de ella; el problema está en quién queda fuera. La cobertura es
de 13.29\,\% de la \zmvm{}, con 53.52\,\% en la \cdmx{} y 3.92\,\% en los
municipios conurbados, 42 de los cuales no tienen ninguna estación a distancia
caminable (§\ref{subsec:cobertura-espacial}). La mitad de los nodos tiene siete
conexiones o menos. El tiempo de viaje promedio es de 113.91~min, y la
intermediación se concentra en Tacubaya ($B(v) = 0.2232$) y en un corredor que
va de Garibaldi a Barranca del Muerto, con un Gini de 0.8007
(§\ref{subsec:fase3-centralidad}). La red, en suma, no tiene un problema de
operación, sino de distribución.

Esa asimetría tiene una lectura social. La \zmvm{} se puede ordenar en tres
contornos de acceso al transporte; en el tercero, que es la corona periférica,
llegar al Metro exige dos o tres transbordos en Corredores Concesionados,
\textsc{rtp} o servicios no formalizados, de modo que el viaje cuesta tiempo y
dinero a la vez. \citet{rendon2009ciudadanos} habla de ciudadanos de primera,
segunda y tercera clase. La clase la define el transporte más cercano. Para
esos habitantes el traslado es también una forma de \termino{pobreza de tiempo}
\citep{gaceta_cuidados_2026}, porque el viaje se come el margen que queda para
el descanso y el cuidado. Sobre la \zmvm{}, entonces, no cabe un juicio único:
una sola red resulta eficiente para una parte del territorio y deficiente para
otra.

El escenario del anillo se lee sobre esa línea base. Su primera lección es que
acercar no es conectar (§\ref{subsec:anillo-conclusiones-contornos}). La
cobertura apenas pasa de 13.29 a 13.43\,\% en los dos casos y los municipios
sin estación cercana bajan de 42 a 41, porque el Periférico corre por la parte
de la ciudad que ya tiene más servicio. Lo que sí se mueve es la conexión de
quien ya está cerca de la red: el tiempo de viaje promedio baja 4.1\,\% con el
Anillo MB y 10.1\,\% con el Anillo Metro (§\ref{subsec:anillo-tiempo-viaje}).
Nezahualcóyotl, Chimalhuacán y Ecatepec ganan entre media hora y una hora con
cualquiera de los dos. El sur depende del modo. Tlalpan, Xochimilco, Tláhuac y
La Magdalena Contreras ganan seis o siete minutos con el primero y entre 15 y
27 con el segundo.

La segunda lección es que el anillo cambia la forma de la red, pero no reparte
su carga (§\ref{subsec:anillo-conclusiones-centralidad}). Con el Anillo MB,
Tacubaya pierde 3.1\,\% de su intermediación y los diez nodos principales siguen
en el centro. Con el Anillo Metro pierde 29.7\,\%, y nueve de esos diez lugares
pasan a estaciones del Periférico. La red deja de ser solo radial y se vuelve
anillo-radial, con un segundo eje sobre el arco sur y oriente. El Gini, sin
embargo, sube a 0.8134: la concentración cambia de lugar, no se reparte. La
diferencia entre los casos la marca el modo
(§\ref{subsec:anillo-conclusiones-modo}). A unos 14~km/h, el Anillo MB ocupa el
papel de la \textsc{rtp} que hoy circula por el Periférico; a 36~km/h, el Anillo
Metro compite con las líneas radiales y reordena la jerarquía. Por eso el
Caso~3 es el más favorable.

El Caso~3 deja además un tipo de nodo que la red actual no tiene. Seis
estaciones del anillo, entre ellas Constitución de 1917, Canal Nacional y
Barranca del Muerto, alcanzan una intermediación de 0.12 a 0.16 con solo 6 a 12
conexiones: por ellas pasarían muchos caminos, pero llegan pocas rutas
(§\ref{subsec:anillo-conclusiones-cetram}). La recomendación es planearlas como
\textsc{cetram} de nueva generación, con rutas alimentadoras y accesos de
última milla. Hay también una lectura histórica. El Periférico se construyó para
el automóvil en los años en que el Metro se rechazaba por costoso
(§\ref{subsec:sesgo-vial} y §\ref{subsec:rechazo-metro}), y en 1985 sus dos
corredores salieron del Programa Maestro por captación. Los resultados sugieren
que esa vialidad es la que mejor podría alojar un transporte masivo lateral, si
tiene vía exclusiva. Eso era lo que no se veía cuando solo se medía la demanda.

El aporte urbano de la tesis es, primero, un diagnóstico de la red
metropolitana completa, con la \cdmx{} y el Estado de México juntos, que queda
como línea base pública y reproducible. Segundo, precisa lo que suele decirse
en general sobre la periferia: lo que falta ahí es cobertura y capilaridad más
que velocidad, y la dependencia de la red tiene nombres y lugares. Tercero,
muestra que el modo y el derecho de vía cambian el tipo de efecto de un
corredor, no solo su tamaño. Quedan preguntas abiertas. El modelo no dice
cuántas personas viven cerca de cada estación, cuánto costaría cada opción, qué
condiciones físicas tiene el trazo ni si dos ejes harían a la red menos frágil
(§\ref{subsec:corredor-anillar-limitaciones}). Tampoco resuelve un problema de
escala. Con el crecimiento disperso de la \zmvm{}, el Periférico ya no es la
periferia: hoy marca el borde exterior del Primer y el Segundo Contorno
(Figura~\ref{fig:semovi-contornos}), y la periferia actual pediría un anillo
exterior. Lo que el modelo no mide define la agenda de
§\ref{sec:conclusiones-academico}.

% !TeX root = ../tesis.tex
%%%%%%%%%%%%%%%%%%%%%%%%%%%%%%%%%%%%%%%%%%%%%%%%%%%%%%%%%%%%%%%%%%%
% SECCIÓN 7.3 — EJE ACADÉMICO Y TRABAJO FUTURO (Issue #23)
% Papel en el hilo: lo que falta medir, y cómo unir estructura y
% demanda. Esquema: aporte Garibelt → L1 consolidar Garibelt →
% L2 otras redes (país y extranjero) → L3 calibrar lo medido →
% L4 medir lo que falta (doctorado) → cierre con 1985.
% Plan: CONTEXTO_CAP7.md §6. Decisiones 2026-11-04: sin IGISC 2026;
% sin nombrar ciudades concretas en L2 (sería especulativo).
%%%%%%%%%%%%%%%%%%%%%%%%%%%%%%%%%%%%%%%%%%%%%%%%%%%%%%%%%%%%%%%%%%%
\section{Eje Académico y Trabajo Futuro}
\label{sec:conclusiones-academico}

El tercer eje reúne lo que la tesis aporta al método y lo que deja abierto. Su
aporte principal es la Clasificación Garibelt, con el Espectro MJG como
entregable (§\ref{subsec:clasificacion-garibelt}). Una red tan heterogénea como
la de la \zmvm{}, con doce sistemas, 76 demarcaciones y una cobertura que cae
de forma abrupta al salir de la \cdmx{} (§\ref{sec:conclusiones-urbano}), no se
deja resumir en un solo número. Un promedio de indicadores esconde la dimensión que está mal; es
el sesgo de compensabilidad que señalan \citet{hossain2024} y \citet{deng2023}
(§\ref{subsubsec:garibelt-compensabilidad}). La Clasificación normaliza para
calibrar, no para sumar, como el perfil por ciudad de \citet{kujala2018}, y
asigna además a cada nodo un intervalo dominante. Así, la red actual sirvió de
línea base y los escenarios se compararon contra ella: diez nodos cambian de
intervalo dominante en el Caso~2 y 21 en el Caso~3 (§\ref{subsec:anillo-espectro-nodo}).

Se trata, con todo, de un instrumento en desarrollo, y la primera línea de
trabajo futuro es consolidarlo. Sus cortes y cotas son supuestos. Bajar la cota
de $DI$ de 2.5 a 2.0 basta para que el $DI$ de la red actual cambie de intervalo
(§\ref{subsubsec:garibelt-escala}), y los valores de $T$ y $B(v)$ dependen de
los parámetros $\alpha$ y $\beta$ del coeficiente de fricción
(§\ref{subsec:fase2-cf}). Hace falta un análisis de sensibilidad de todos
ellos. Hay también dos medidas relativas: el intervalo dominante de cada nodo se
calcula contra el máximo de cada caso, y la $FC$ se normaliza con los cuartiles
de la propia red, de modo que describe cómo se reparte la conectividad dentro
de la \zmvm{} sin decir cómo se compara con otras
(§\ref{subsubsec:garibelt-procedimiento}). Resueltos esos puntos, la
Clasificación podría servir no solo para diagnosticar una red heterogénea, sino
para ordenar varias entre sí.

Ahí empieza la segunda línea. La tesis midió una sola red real frente a dos
casos hipotéticos construidos sobre ella, así que sus resultados solo se
comparan consigo mismos. El paso siguiente es aplicar la cadena a otras redes:
las de otras zonas metropolitanas del país y las de ciudades de otros países
que tengan \termino{feeds} \gtfs{} y cartografía de base
(§\ref{subsec:datos-herramienta-replicable}). Con varios perfiles del Espectro
MJG, como los 25 que reunieron \citet{kujala2018}, se sabría si los valores de
la \zmvm{} son altos o bajos frente a otras ciudades. También se podrían
sustituir los cuartiles propios de $FC$ por una referencia común. Y la
replicabilidad que planteó §\ref{sec:conclusiones-tecnologico} quedaría
probada fuera de la ciudad donde se construyó la herramienta.

La tercera línea es calibrar lo que el modelo ya mide. La cobertura podría
calcularse con isócronas peatonales sobre la red vial de OpenStreetMap y
cruzarse con la población por manzana del Censo 2020, para pasar de superficie
cubierta a personas cubiertas (§\ref{subsec:indicador-cobertura}). La
penalización por transbordo podría contrastarse con los tiempos que devuelve la
Google Maps Directions API en horas pico y valle
(§\ref{subsec:penalizacion-transferencia}). Las velocidades de superficie
tendrían que compararse con velocidades observadas por ruta. La $FC$, por su
parte, podría ponderarse por capacidad o por flujo, como en la definición de
\citet{lee2008}, y contrastarse con el catálogo oficial de \textsc{cetram}
(§\ref{subsec:fuerza-capilar}).

La cuarta línea es medir lo que hoy queda fuera, y es la que da forma a una
propuesta de nivel doctoral. La robustez $\Delta E$ pide la eficiencia global
de \citet{latora2001efficient}, la remoción de Macro-Hubs completos y una
corrida vigente (§\ref{subsec:vulnerabilidad-estructural}); solo así se sabría
si la red anillo-radial resiste mejor una falla en el centro. A eso se suman la
demanda, la población, el ingreso y el costo de cada opción, que el modelo no
incorpora (§\ref{subsec:corredor-anillar-limitaciones}), y un Anillo Periférico
Exterior para la periferia actual. El núcleo de esa agenda es unir en un mismo
modelo los indicadores de estructura con la afluencia y las matrices
origen-destino.

Esa unión es también el cierre del argumento. En 1985, el Programa Maestro
preguntó cuánta gente usaría un corredor sobre el Periférico y, con esa sola
medida, lo descartó. Esta tesis preguntó qué aportaba ese corredor a la
estructura de la red y encontró efectos que la captación no registraba. Ninguna
de las dos preguntas basta por separado. La planeación por estructura no
sustituye a la planeación por demanda; la completa. Si lo que se mide define lo
que se construye, el trabajo que sigue es medir las dos cosas a la vez.

